\section{Equilibrium bifurcation structure}
\label{sec:results}

We now apply the numerical framework of Section~\ref{sec:methods} to determine how the
restricted equilibria change as the underlying four-primary central
configuration is varied. We first establish the four reference
configurations used as the computational baseline (Section~\ref{subsec:verification}). We then
examine the common Case-1/Case-2 slice and determine the component structure
of its admissible central configurations (Section~\ref{subsec:components}), map the equilibrium
multiplicity along the resulting admissible arcs and resolve the
eleven-equilibrium regime (Section~\ref{subsec:cascades}), and finally verify every interior
multiplicity change as a generic saddle-node fold (Section~\ref{subsec:foldverify}).


\subsection{Baseline configurations}
\label{subsec:verification}

The four reference configurations are first refined onto the
central-configuration manifold using the procedure of Section~\ref{subsec:eqsearch}. At each
refined configuration the mass ratios are recomputed from
\eqref{eq:mass_ratios}, admissibility is checked, and the restricted
equilibria are located independently by the multi-start search. The results
are summarized in Table~\ref{tab:baseline}. Every retained equilibrium
satisfies
\[
\|\nabla\Omega\|<10^{-12}.
\]

The resulting equilibrium counts are
\begin{equation}
N_{\rm eq}=5,\;9,\;5,\;5
\qquad
\text{for Cases~1--4, respectively}.
\label{eq:baseline_counts}
\end{equation}
Thus the finite-precision reference geometries can be placed consistently
on the central-configuration manifold without altering their equilibrium
multiplicities.

One feature of this baseline will become important in Section~\ref{sec:stability}. Case~1
contains an outer equilibrium near
\[
(x,y)\simeq(2.67,\,2.62)
\]
whose linear spectrum consists of two purely imaginary conjugate pairs. It
is therefore spectrally stable, whereas the remaining equilibria of the
four reference configurations are spectrally unstable. The complete
equilibrium coordinates, refinement displacements, and per-root numerical
diagnostics are provided in the accompanying reproducibility archive.

\begin{table}[t]
\centering
\small
\caption{The four representative non-symmetric configurations used as the
computational baseline. The approximate geometry is refined onto the
central-configuration manifold $\psi=0$ while holding $(a,b)$ fixed, after
which the mass ratios are recomputed. $N_{\rm eq}$ is the equilibrium count
returned by the independent multi-start search, with every retained root
satisfying $\|\nabla\Omega\|<10^{-12}$.}
\label{tab:baseline}
\begin{tabular}{lccccc}
\hline
Case &
$(a,b)$ &
$(c_1,c_2)_{\rm approx}$ &
$(c_1,c_2)_{\rm refined}$ &
$(\rho_0,\rho_1,\rho_2)$ &
$N_{\rm eq}$\\
\hline
Case 1 &
$(0.9,-0.1)$ &
$(-0.61,0.71)$ &
$(-0.610000,0.715655)$ &
$(49.654,3.0981,1.1533)$ &
5\\

Case 2 &
$(0.9,-0.1)$ &
$(0.52,-0.88)$ &
$(0.517714,-0.880000)$ &
$(4.2776,0.95681,1.2772)$ &
9\\

Case 3 &
$(1.5,0.4)$ &
$(-0.96,0.99999)$ &
$(-0.944942,0.999990)$ &
$(3.0082,0.15809,0.43247)$ &
5\\

Case 4 &
$(1.5,0.4)$ &
$(1.05,0.13)$ &
$(1.050584,0.130000)$ &
$(4.3119,0.98292,0.72214)$ &
5\\
\hline
\end{tabular}
\end{table}


\subsection{Distinct central-configuration components on the fixed slice}
\label{subsec:components}

Cases~1 and~2 share the same values
\[
(a,b)=(0.9,-0.1),
\]
while possessing different equilibrium multiplicities,
$N_{\rm eq}=5$ and $N_{\rm eq}=9$, respectively. Their common
fixed-$(a,b)$ slice therefore provides a natural controlled setting in
which to ask whether the two configurations belong to the same admissible
central-configuration family and how the restricted equilibrium structure
changes along that family. The slice is selected directly from the
reference configurations, before the continuation analysis; the
disconnected component structure and the multiplicity cascade reported
below are outcomes of the continuation rather than criteria used to choose
the slice.

We consequently continue
\begin{equation}
\psi(0.9,-0.1,c_1,c_2)=0
\label{eq:fixed_slice_constraint}
\end{equation}
by the pseudo-arclength method of Section~\ref{subsec:continuation}. The continuation reveals two
qualitatively different zero-set branches containing the reference
configurations: an open branch through Case~1 and a closed loop through
Case~2 (Figure~\ref{fig:cc_arcs}). Restricting these branches to positive
mass ratios produces the two admissible arcs used throughout the remainder
of the paper.

The closed character of the Case-2 branch is obtained directly from the
pseudo-arclength trace: continuation around the branch in the
$(c_1,c_2)$ plane returns to the starting neighbourhood and completes the
loop. No vanishing of
\[
\nabla_{(c_1,c_2)}\psi
\]
is encountered along the accepted continuation points. Thus, within the
resolution of the numerical trace, the Case-2 loop is a regular level-set
component. At a regular point the implicit-function theorem gives a locally
unique smooth zero-set curve, so another branch cannot join or leave the
loop there. The Case-1 reference configuration lies instead on the
separately traced open branch.

The numerical continuation therefore supports the conclusion that
Cases~1 and~2 lie on distinct connected components of the regular level set
$\psi=0$ on the investigated fixed-$(a,b)$ slice. This is deliberately a
slice-specific numerical statement, not a global theorem for the complete
four-parameter central-configuration set. In particular, the present
calculation does not exclude the possibility that a path connecting the two
configurations exists when $a$ and $b$ are also allowed to vary.

The consequence for the restricted problem is important. The five- and
nine-equilibrium reference configurations are not two points on a single
fixed-$(a,b)$ continuation family. Their different multiplicities therefore
cannot be interpreted as evidence of a bifurcation sequence directly
joining Case~1 to Case~2. Equilibrium multiplicity must instead be analysed
separately on the two components.

On the branch through Case~1, the maximal admissible arc containing the
reference configuration extends from a boundary where
$\rho_0\rightarrow0$ to one where $\rho_2\rightarrow0$. On the branch
through Case~2, the corresponding admissible arc extends from
$\rho_1\rightarrow0$ to $\rho_0\rightarrow0$. Each maximal admissible arc
is parametrized below by normalized arclength
\[
s\in[0,1].
\]

\begin{figure}[htbp]
\centering
\includegraphics[width=0.72\textwidth]{cc_continuation_path.pdf}
\caption{Central-configuration branches and admissible arcs in the
$(c_1,c_2)$ plane at $(a,b)=(0.9,-0.1)$. Grey shows the numerically
represented $\psi=0$ branches in the plotted domain; blue and red denote
the admissible arcs containing Cases~1 and~2, respectively. Case~1 is
marked by a star and Case~2 by a diamond. Gold circles mark the seven
verified interior saddle-node folds $F_1$--$F_7$; the inset resolves the
closely spaced folds $F_4$--$F_7$. The two reference configurations lie
on distinct fixed-slice components, and each admissible arc terminates
where a primary mass ratio vanishes.}
\label{fig:cc_arcs}
\end{figure}


\subsection{Fold cascades and the eleven-equilibrium regime}
\label{subsec:cascades}

The equilibrium multiplicity cannot be determined by continuing only the
branches already known at a reference configuration, because a pair born
at a saddle-node does not exist on the opposite side of its fold. We
therefore use the independent multi-start search described in
Section~\ref{subsec:continuation} at every accepted central-configuration sample.

The resulting multiplicity function $N_{\rm eq}(s)$ is shown in
Figure~\ref{fig:count_staircase}. Along the two admissible arcs we obtain
\begin{equation}
\begin{aligned}
\text{Case-1 arc:}&\qquad
5\rightarrow7\rightarrow5\rightarrow4,
\\
\text{Case-2 arc:}&\qquad
7\rightarrow9\rightarrow11\rightarrow9\rightarrow7\rightarrow5.
\end{aligned}
\label{eq:count_cascades}
\end{equation}

The seven interior changes all have magnitude two and are subsequently
verified as generic saddle-node folds. They are denoted
$F_1,\ldots,F_7$ in Figures~\ref{fig:cc_arcs}--\ref{fig:count_staircase}
and listed in Table~\ref{tab:bifurcations}.

The final $5\rightarrow4$ change on the Case-1 arc is qualitatively
different. It occurs near
\[
s\simeq0.974
\]
as $\rho_2\rightarrow0$, when the underlying positive-mass primary
configuration approaches the boundary of the admissible set and an
associated equilibrium recedes to large distance. This is therefore a
mass-vanishing parameter-boundary degeneration, not an interior
saddle-node bifurcation. More generally, the admissible CC arcs terminate
when one of the primary mass ratios vanishes, but such arc termination must
not be conflated with an interior bifurcation of the restricted equilibrium
equations.

The local primary searches of Section~\ref{subsec:continuation} are especially important near
these limits, where equilibria associated with a small primary mass can lie
close to the corresponding singularity. The multiplicities in
\eqref{eq:count_cascades} are those obtained after applying the augmented
global-plus-local search throughout the arcs.

\begin{figure}[htbp]
\centering
\includegraphics[width=\textwidth]{equilibrium_count_vs_s.pdf}
\caption{Equilibrium multiplicity $N_{\rm eq}(s)$ along the admissible
arcs. Top: Case-1 arc; bottom: Case-2 arc. Dashed vertical lines mark the
verified saddle-node folds $F_1$--$F_7$. The right-most Case-1 transition
is the mass-vanishing event $\rho_2\rightarrow0$ and is not an interior
fold.}
\label{fig:count_staircase}
\end{figure}


\begin{table}[t]
\centering
\footnotesize
\setlength{\tabcolsep}{4pt}
\caption{Verified changes in equilibrium multiplicity along the two
admissible central-configuration arcs. Interior count changes are generic
saddle-node folds located from
$\Omega_x=\Omega_y=\det H_\Omega=\psi=0$ and satisfying the
nondegeneracy conditions of Section~\ref{subsec:foldmethods}. The coordinates
$(c_1,c_2)_c$ describe the critical \emph{primary configuration};
$(\alpha,\beta)$ are the transversality and quadratic-nondegeneracy
coefficients of \eqref{eq:fold_coefficients}. The Case-1
$5\rightarrow4$ event is a mass-vanishing boundary degeneration rather
than an interior fold.}
\label{tab:bifurcations}
\begin{tabular}{llcccccl}
\hline
Arc & $N_{\rm eq}$ & Fold & $s_c$ &
$(c_1,c_2)_c$ & $(\rho_0,\rho_1,\rho_2)_c$ &
$(\alpha,\beta)$ & Type\\
\hline

Case-1 &
$5\to7$ &
$F_1$ &
0.261 &
$(-1.0717,1.0890)$ &
$(0.507,0.993,0.965)$ &
$(-2.72,22.4)$ &
saddle-node\\

Case-1 &
$7\to5$ &
$F_2$ &
0.583 &
$(-0.8206,0.8193)$ &
$(3.95,1.78,1.65)$ &
$(16.3,42.6)$ &
saddle-node\\

Case-1 &
$5\to4$ &
n/a &
0.974 &
n/a &
n/a &
n/a &
$\rho_2\to0$ boundary\\

Case-2 &
$7\to9$ &
$F_3$ &
0.366 &
$(0.4040,-0.9842)$ &
$(8.9,0.667,1.27)$ &
$(-0.615,0.0566)$ &
saddle-node\\

Case-2 &
$9\to11$ &
$F_4$ &
0.768 &
$(0.4478,-0.6231)$ &
$(0.71,2.19,2.28)$ &
$(-18.8,43.4)$ &
saddle-node\\

Case-2 &
$11\to9$ &
$F_5$ &
0.804 &
$(0.4246,-0.5923)$ &
$(0.576,2.49,2.57)$ &
$(1.18,0.231)$ &
saddle-node\\

Case-2 &
$9\to7$ &
$F_6$ &
0.839 &
$(0.3998,-0.5606)$ &
$(0.456,2.87,2.94)$ &
$(-6.03,-48.8)$ &
saddle-node\\

Case-2 &
$7\to5$ &
$F_7$ &
0.929 &
$(0.3333,-0.4795)$ &
$(0.189,4.3,4.35)$ &
$(0.506,0.173)$ &
saddle-node\\

\hline
\end{tabular}
\end{table}


\subsubsection{The eleven-equilibrium regime}
\label{subsubsec:eleven}

The most pronounced multiplicity increase occurs on the Case-2 arc. The
$9\rightarrow11$ fold $F_4$ at
\[
s_c\simeq0.768
\]
and the $11\rightarrow9$ fold $F_5$ at
\[
s_c\simeq0.804
\]
bound a narrow interval in which eleven distinct restricted equilibria are
detected. The two folds occur at different primary configurations
(Table~\ref{tab:bifurcations}) and at spatially separated restricted
equilibria. They are therefore two distinct saddle-node events rather than
a single higher-codimension degeneracy.

Tracking the individual branches through this interval reveals a feature
that is not visible from the multiplicity alone. At $F_4$, a new pair,
denoted $B_a$, is created near the critical restricted equilibrium
\[
(x,y)\simeq(-0.38,\,0.31).
\]
At the representative point
\[
s\simeq0.786
\]
inside the eleven-equilibrium interval, its two members lie approximately
at
\[
(-0.439,\,0.377),
\qquad
(-0.301,\,0.218).
\]

The pair destroyed at the subsequent $11\rightarrow9$ fold $F_5$ is not
$B_a$. Instead, $F_5$ annihilates a pre-existing pair, denoted $B_b$, in
the lower part of the equilibrium configuration near
\[
(x,y)\simeq(0.53,\,-1.70).
\]
Immediately before the fold its two members lie approximately at
\[
(0.26,\,-1.74),
\qquad
(0.74,\,-1.64).
\]
Branch tracking shows that this pair was created earlier at the
$7\rightarrow9$ fold $F_3$.

The newly created pair $B_a$ therefore survives the
$11\rightarrow9$ transition and remains present throughout the following
nine-equilibrium interval. It is annihilated only at $F_6$, the subsequent
$9\rightarrow7$ fold near
\[
(x,y)\simeq(-0.26,\,0.12).
\]
Consequently,
\[
9\rightarrow11\rightarrow9
\]
is not a simple creation and subsequent destruction of the same equilibrium
pair. The two nine-equilibrium intervals on either side of the
eleven-equilibrium window have the same multiplicity but different branch
populations. The eleven-equilibrium interval therefore exhibits a genuine
branch exchange.

As a robustness check, the equilibrium search was repeated at the
representative point
\[
s\simeq0.786,
\qquad
(c_1,c_2)\simeq(0.432,-0.602),
\]
inside the eleven-equilibrium interval.
The re-search uses the same $81\times81$ deterministic grid over
$[-8,8]^2$ and the same $35\times35$ local grid around each primary as the
main continuation analysis, but increases the number of fixed-seed random
starts from 1000 to 1500 (seed 20260808). It again returns eleven distinct
equilibria. Every retained root satisfies
\[
\|\nabla\Omega\|<10^{-12}.
\]
The minimum pairwise separation among all eleven roots at this configuration
is approximately
\[
0.175,
\]
while the separation of the two members of the newly created pair $B_a$ is
approximately
\[
0.21.
\]

We therefore report eleven as the number of distinct equilibria detected
under the stated numerical search procedure. This is a numerically supported
multiplicity on the investigated parameter interval; it is not an analytic
upper bound on the number of real equilibria of the general restricted
five-body problem.

This count is stable under refinement of the search. Repeating the multi-start
search at the eleven-equilibrium configuration over grid densities $61^2$
through $121^2$, random-start counts from $1000$ to $4000$, and three
independent seeds returns eleven equilibria in every case. Reducing the
converged-root rejection radius from $10^{-2}$ through $5\times10^{-3}$ to
$10^{-3}$ reveals no additional roots, and the closest equilibrium to any
primary at this configuration lies at distance $0.33$, so no genuine
near-primary equilibrium is discarded by the exclusion guard. The same
invariance holds for the four reference configurations, whose nearest
equilibrium-to-primary distance is $0.13$. These convergence and
exclusion-radius data are provided in the reproducibility archive.

Figure~\ref{fig:case2_cascade} shows representative equilibrium geometries
through the complete Case-2 sequence
\[
7\rightarrow9\rightarrow11\rightarrow9\rightarrow7\rightarrow5,
\]
with $B_a$ and $B_b$ identified explicitly. The radial branch representation
in Figure~\ref{fig:radial_branches} gives the complementary continuation
view and makes the branch exchange visible even where the scalar
multiplicity alone cannot distinguish branch identity.

\begin{figure}[htbp]
\centering
\includegraphics[width=\textwidth]{cascade_geometry_snapshots.pdf}
\caption{Equilibrium geometry at representative points of the Case-2
cascade $7\rightarrow9\rightarrow11\rightarrow9\rightarrow7
\rightarrow5$. Stars denote primaries, red circles unstable equilibria,
and green diamonds spectrally stable equilibria. The branch exchange in
the eleven-equilibrium interval is annotated explicitly: $B_a$, created
at $F_4$, persists into the following $N_{\rm eq}=9$ interval, whereas
the pre-existing pair $B_b$ is destroyed at $F_5$.}
\label{fig:case2_cascade}
\end{figure}


The branch identities shown in Figure~\ref{fig:radial_branches} are those
obtained from the adaptive predictor-based tracker of Section~\ref{subsec:continuation}. The same
branch assignment is used for the stability, centre-frequency, Jacobi, and
periodic-orbit analyses that follow. Ordinary equilibrium branches are
therefore shown without individual labels, while the dynamically important
pairs $B_a$ and $B_b$ are highlighted explicitly. Large-radius excursions
near the ends of the admissible arcs correspond to equilibria receding as a
primary mass ratio approaches zero.

\begin{figure}[htbp]
\centering
\includegraphics[width=\textwidth]{equilibrium_branches_r_vs_s.pdf}
\caption{Radial equilibrium coordinate
$r_i(s)=\sqrt{x_i(s)^2+y_i(s)^2}$ along the two admissible arcs.
Ordinary branches are shown in muted grey and dashed vertical lines mark
folds $F_1$--$F_7$. Spectrally stable equilibria are highlighted as
described in Section~\ref{sec:stability}. On the Case-2 arc, $B_a$ is the pair created at
$F_4$, whereas $B_b$ is the pre-existing pair destroyed at $F_5$,
directly displaying the branch exchange in the eleven-equilibrium
interval. Large-radius excursions near the arc ends correspond to
equilibria receding at mass-vanishing limits.}
\label{fig:radial_branches}
\end{figure}


\subsection{Saddle-node verification and scaling}
\label{subsec:foldverify}

The seven interior multiplicity changes in
\eqref{eq:count_cascades} are tested using the fold criteria of
Section~\ref{subsec:foldmethods}. For every transition, the augmented system
\begin{equation}
\Omega_x=\Omega_y=\det H_\Omega=\psi=0
\label{eq:fold_conditions_results}
\end{equation}
converges to a distinct critical point with augmented-system residual not
exceeding $10^{-11}$.

At each critical point the Hessian has one numerically vanishing eigenvalue,
denoted $\mu_0$, and one nonzero transverse eigenvalue, denoted
$\mu_\perp$. The corresponding transversality and
quadratic-nondegeneracy coefficients satisfy
\[
\alpha\neq0,
\qquad
\beta\neq0
\]
for all seven folds, with the computed $(\alpha,\beta)$ values listed in
Table~\ref{tab:bifurcations}. Thus every interior count change satisfies the
numerical conditions for a simple generic saddle-node fold.

The local branch geometry provides an independent verification. For each
fold the merging pair is resolved on the side of the critical parameter on
which both equilibria exist. The two roots approach one another and meet at
the numerically determined fold location, as illustrated for representative
cases in Figure~\ref{fig:fold_local}.

We further fit their separation to
\begin{equation}
d\propto|\ell-\ell_c|^\gamma
\label{eq:measured_fold_scaling}
\end{equation}
using the innermost resolved points. Across the seven folds the fitted
exponents satisfy
\[
0.44\leq\gamma\leq0.52,
\]
consistent with the Lyapunov--Schmidt prediction
\[
\gamma=\frac{1}{2}
\]
derived in Section~\ref{subsec:foldmethods}. Since the normalized coordinate $s$ is a smooth
monotone rescaling of $\ell$ on each admissible arc, the same local exponent
is obtained when the distance to the critical point is expressed using
$|s-s_c|$.

The agreement of four independent diagnostics (the observed
$\Delta N_{\rm eq}=\pm2$ count changes, convergence of the augmented
degeneracy system, satisfaction of the fold nondegeneracy conditions, and
the square-root pair-separation law) provides mutually consistent evidence
that $F_1,\ldots,F_7$ are generic saddle-node bifurcations of the restricted
equilibrium equations.

\begin{figure}[htbp]
\centering
\includegraphics[width=\textwidth]{fold_zoom_main.pdf}
\caption{Three representative saddle-node folds (a Case-1 fold, an
ordinary Case-2 fold, and the $9\rightarrow11$ fold bounding the
eleven-equilibrium interval): the $x$-coordinates of the merging pair meet
at $\ell=\ell_c$ (dashed).}
\label{fig:fold_local}
\end{figure}

\begin{figure}[htbp]
\centering
\includegraphics[width=0.78\textwidth]{fold_separation_exponent.pdf}
\caption{Saddle-node pair-separation scaling. Left: a representative fold
(Case-2 $7\rightarrow9$) showing the merging-pair separation $d$ against
$|\ell-\ell_c|$ on log--log axes, with the fitted slope and the reference
slope $\tfrac12$. Right: the fitted exponents $\gamma$ of all seven folds
against the theoretical value $\gamma=\tfrac12$ (shaded band
$[0.45,0.55]$); all lie in $[0.44,0.52]$.}
\label{fig:fold_scaling}
\end{figure}